\chapter*{Conferencia VII}
\addcontentsline{toc}{chapter}{LECTURA I - El timbre y el carácter del violín}
SEÑORES: En este momento les presento el barniz.
Fenómeno n.° 2. Como recordará, el fenómeno del barniz n.° 1 ocurrió en la superficie exterior de mi violín. El fenómeno del barniz n.° 2 ocurrió en la superficie interior de otro violín. Este último violín es el primero de una lista de sesenta en los que apliqué una protección para la superficie interior, y este último fenómeno se debe a dos errores:
(1) Tosquedad en la aplicación del material. (2) Ensamblaje y prueba del tono antes de que el material se vuelva
seco.
Pensando que algún lector podría desear restaurar la superficie interior de un violín usado, y con el fin de evitarle decepciones, me aseguro de señalar mis propios errores. Es lógico suponer que, en cualquier actividad humana, se cometen errores. Al conocerlos, se pueden evitar. Tras diez años de esfuerzo protegiendo la superficie interior del violín de la desintegración y creyendo haberlo logrado sin afectar el sonido, no afirmo que los detalles de mi trabajo no puedan mejorarse. Me alegraría enormemente saber que mis esfuerzos han mejorado.
La tosquedad en la aplicación del material se muestra así: (a) Por falla en la preparación de la superficie de la madera, (b) Por aplicación de un exceso de material. Después de mi experiencia en este trabajo de revestimiento interior, después de observar la intensidad y el brillo del tono sin pérdida alguna de dulzura, me encuentro convencido.
113

Me pregunto por qué no comprendí antes que las superficies interiores del violín deberían ser igual de perfectas que las de la flauta, el clarinete y la trompa. En estos últimos instrumentos, la intensidad y el brillo del sonido se modifican, para bien o para mal, según el grado de perfección de la superficie interior. Ahora estoy completamente convencido de que la misma condición afecta al sonido del violín. En cualquier método de aplicación de un acabado transparente a la madera, la hinchazón de la veta es una dificultad a tener en cuenta. Yendo directamente al grano, no conozco nada que cause menos hinchazón de la veta que el aceite de linaza hervido. Sin embargo, he descubierto que hay una forma correcta e incorrecta de aplicar el aceite. No he encontrado que el aceite solo, o mezclado con gomas de secado más lento, penetre ni siquiera en el pino o cedro blanco más blando cuando se aplica en capas tan finas como las que permite el proceso de frotado. He descubierto que las superficies de madera deben prepararse cuidadosamente para asegurar la mayor atenuación del aceite, ya sea solo o en mezclas. En violines usados, nunca he encontrado superficies interiores listas para recibir el material de acabado. La superficie interior de la caja de resonancia usada suele presentar una serie de crestas y valles causados por la hinchazón de los tejidos conectivos entre las fibras más densas. La profundidad y altura de estas crestas y valles varían con la cantidad de vapor de agua absorbido. Como se ha demostrado anteriormente, diferentes muestras de madera de caja de resonancia absorben diferentes cantidades de humedad debido a variaciones en el calibre de los tubos capilares que transportan la savia. La superficie interior del fondo y las costillas de madera dura a menudo presenta irregularidades.
114

Contracción uniforme en las marcas transversales; la parte más densa se presenta como una cresta baja, la parte tubular como una depresión poco profunda. El objetivo es nivelar todas estas crestas y eliminar toda la humedad antes de aplicar cualquier material protector. A menudo, considero necesario el uso de calor artificial para eliminar la humedad de la madera de violín usada; también es necesario continuar aplicando calor artificial hasta que se complete el trabajo de superficie y la aplicación del material de acabado.
hechos.
Para nivelar las crestas, prefiero usar papel de lija sobre un bloque de madera en lugar de sostenerlo sobre los dedos. Las yemas de los dedos, al ser suaves, amortiguan más o menos, y así continúan eliminando las superficies de los valles, mientras que el bloque, al ser rígido, solo elimina las crestas. Para tal bloque
Utilizo cualquier madera firme, de 1 pulgada de espesor y H pulgadas de largo.
de ancho y 2i de largo, con un borde curvo y ligeramente ovalado. En la caja de resonancia, las crestas se nivelan trabajando a través de las marcas transversales. A veces, este trabajo en la placa posterior requiere bastante tiempo para reducir las densas crestas. Pero, en mis manos, el raspador es una herramienta peligrosa para tal trabajo, debido a que arranca trozos de madera, y especialmente peligroso en madera vieja y quebradiza. Por tales razones, continúo pacientemente con el papel de lija hasta que la superficie marrón de los valles parezca nueva. Luego elimino las marcas del papel de lija con piedra pómez en polvo y una almohadilla de pulido de fieltro con una superficie ligeramente humedecida con aceite hervido, usando solo la cantidad suficiente de aceite.
115

para mantener la piedra pómez en su trabajo. El trabajo preparatorio final se realiza con el "hueso". El "hueso" es un trozo de ébano duro, de un tamaño conveniente para sostener, con un extremo o borde ligeramente curvado y ovalado, y pulido al máximo. Con este "hueso" se frotan las superficies hasta que alcanzan una suavidad similar a la de un espejo, después de lo cual están listas para recibir el material protector. Sobre las superficies así preparadas, se puede aplicar aceite, solo o mezclado, en capas de la mayor atenuación posible, con una hinchazón imperceptible del grano. La superficie interna de las costillas, los bloques y los revestimientos reciben la misma atención cuidadosa.
En este trabajo hay dos objetivos importantes a tener en cuenta. Primero, aplicar solo el material suficiente para cubrir la madera. Segundo: Producir una superficie perfectamente lisa. Para este trabajo, el pincel no puede asegurar la atenuación del material que se logra fácilmente mediante el proceso de frotado. La fluidez del material protector es algo que se debe evitar, especialmente en la primera capa. La madera seca absorberá la humedad del líquido aplicado por el pincel, frustrando así el objetivo de sellar herméticamente la MADERA SECA. El caso es diferente al de aplicar protección a la superficie exterior de las tapas del violín. Si se produce absorción de líquido en la superficie exterior, puede escapar por la superficie interior, y ninguna precaución contra la entrada de humedad en la madera puede ser excesiva. Cuando está sellada herméticamente, la humedad no puede entrar ni salir de la madera. De ahí la necesidad de realizar este trabajo en un ambiente seco y
116

Empleo de un material con el mínimo de líquido. No conozco ningún material protector que absorba menos humedad que el aceite de linaza hervido. Que este aceite no penetra en la madera más blanda, cuando se aplica mediante frotamiento y en capas finas, se demostrará más adelante. Pero, tras observar que cuando se emplea una mezcla de aceite y goma copal, atenuada con gomas más blandas, el aceite, al secarse con menos rapidez y ser expulsado en gran medida de la mezcla, quedando sobre la superficie, desde entonces he empleado esta mezcla para la protección de superficies interiores, excluyendo todos los demás agentes. Una vez expulsado el aceite, he probado dos métodos para su eliminación. Primero: Dejar que permanezca y seque como parte del acabado. Segundo: Retirarlo con un paño suave doce horas después de la aplicación. Las gomas, al estar semisólidas, no se ásperan con la cuidadosa eliminación del aceite. Este último método tiene la ventaja de una mayor suavidad de la superficie y una mayor rapidez de secado. Una vez retirado el aceite de la superficie, la goma se secará en 24 horas. Con el aceite superficial restante, el tiempo necesario para el secado será
Varían en cuanto a la cantidad de aceite, de 4 a 7 días.
Para la aplicación del material protector, los detalles son los siguientes: En mi experiencia, la franela de algodón gruesa, fina y de pelo largo ha demostrado ser el mejor material para la almohadilla de frotamiento. Corto un trozo de 2x4 pulgadas y lo doblo una vez, de extremo a extremo, y lo aprieto hacia afuera. Para asegurar el uso de una pequeña cantidad de barniz y aceite a la vez, recurro a los siguientes medios: Un frasco de 2 oz. se llena dos veces.
117

tercios con barniz; un frasco de 2 oz. se llena hasta dos tercios con aceite hervido. Humedezca ligeramente una superficie de la almohadilla con aceite, (sin permitir que se acumule aceite ni barniz), luego sostenga el centro de la almohadilla sobre la boca del frasco de barniz, invierta este último, luego frote la almohadilla sobre la madera con un movimiento circular para evitar que se pegue y la consiguiente aspereza de la superficie a barnizar. La almohadilla debe sujetarse firmemente y con la punta del dedo índice presionando hacia abajo sobre su centro. Cuando esta cantidad de material se haya extendido hasta sus límites, vuelva a los frascos para obtener más aceite y barniz, continuando así hasta que la superficie de la placa esté cubierta. La cantidad de material utilizada para una capa puede parecer sorprendentemente pequeña. Es pequeña. Debe ser pequeña. También debe ser pequeña en la superficie exterior del violín. En cualquier caso, solo uso el material suficiente para cubrir la madera.
Hasta donde sé, este método de aplicación requiere la mínima cantidad de material y logra la máxima perfección superficial. Una vez aplicada con cuidado y seca la última capa, se procede al pulido final con piedra pómez seca en polvo, utilizando una almohadilla nueva de franela de algodón. Por supuesto, se requiere experiencia para dominar este trabajo.
He colocado los detalles correctos antes que los erróneos para que estos últimos resalten con colores llamativos. La tosquedad en la aplicación del material se debe al fracaso de mi primer intento de proteger las superficies interiores del violín; y a la misma causa se debe el fenómeno del barniz n.° 2 que ahora se explicará.
aparecer. 118

El fenómeno del barniz n.° 2 no contribuye en nada al valor del violín ni resulta de interés para el estudiante de violín, más allá de la descripción de detalles erróneos en la protección de la superficie interior. Conociendo el prejuicio existente hacia cualquier tipo de protección para las superficies interiores del violín, me aseguro de describir minuciosamente tanto los métodos correctos como los incorrectos de aplicación. Tras diez años de experiencia en este campo, estoy plenamente convencido de que el prejuicio contra la protección de la superficie interior del violín se debe a detalles erróneos en la aplicación del material protector. Recuerdo haber oído que, como argumento en contra de la protección, esta "aguda" el sonido. Ahora presento pruebas de que dicha protección puede aplicarse de forma que produzca un sonido "apagado".
[Más que el violín, no conozco ningún tema histórico que ofrezca mayor confusión en cuanto a la evidencia. Por alguna razón oculta, el violín no suena ni parece igual para dos personas distintas. Este hecho es un fenómeno sin explicación.]
Calificado por su valor tonal, el violín en el que ocurrió el fenómeno de barniz n.° 2 valía 25 dólares. Este violín había sido usado durante cinco años cuando lo seleccioné para la aplicación de protección de la superficie interior. Para acceder a sus superficies interiores, se retiró la caja de resonancia. No se realizó ninguna preparación más allá de cepillar el polvo de esas superficies. Sobre ellas vertí una cantidad de aceite hervido y barniz de copal transparente templado. Extendí estos materiales con una fregona. Tal como está
119

Bien conocido en la física de la limpieza, el trapeador se hace actuar mediante movimiento recíproco; el movimiento circular se emplea solo en partes que siguen el piso." Sin tener que lidiar con esta última dificultad, puse mi trapeador en acción mediante movimiento recíproco. Por supuesto, el trapeador "se atascaba" en cada movimiento inverso, mientras yo "me aferraba" a mi trabajo vertiendo más aceite y barniz. Después de "mezclarlo" así hasta que la masa parecía como si estuviera "extendida" (¡oh!), declaré que esas superficies estaban ampliamente protegidas. Como mi pegamento estaba caliente, el trabajo de ensamblaje se completó rápidamente. Dos horas después, los gérmenes de la fiebre de "regreso rápido" me habían vuelto loco. Esa sensación de "no puedo esperar" pronto se convirtió en dominante. Apliqué el arco.
[La cabalgata de Sheridan hacia Winchester, donde transformó la derrota en victoria, ha sido retratada por hábiles historiadores; siete poetas han dedicado sus liras al tatuaje de las patas de su caballo. La fama de Sheridan se extendió por todo el mundo en pocas horas. En mi caso, transcurrieron varios años antes de que me atreviera a mirar a la cara a un historiador o a un poeta.] No puedo decir
Mucho por el tono de este violín. A decir verdad, no tiene ningún tono digno de mención. Su tono no solo está "muerto", sino que está "enterrado". El único valor de su tono es un ejemplo de cómo no se debe trabajar la superficie interior y como advertencia contra la fiebre de retorno rápido. Aplicar el arco antes de que el material protector sobre cualquiera de las superficies de las tapas del violín se seque es un acto de locura. Si el historiador hubiera tenido acceso a los hechos "pegajosos" que se adhieren al interior de este violín, entonces el estudiante de violín de hoy 120

Sería como leer que la protección de la superficie interior del violín es un fracaso. Si yo mismo hubiera abandonado el intento después de este desastre, ahora creería que dicha protección causa un tono "muerto". De hecho, lo creí durante los siguientes seis meses. Entonces, en un momento de ocio, tomé este violín desechado y volví a usar el arco. Me esperaba una sorpresa. En cuanto a intensidad y brillantez, su sonido ahora superaba mi valoración anterior. Rápidamente volví a quitar la caja de resonancia.
Quienes hayan leído la descripción del fenómeno del barniz n.° 1 quizás se sientan preparados para la descripción del fenómeno n.° 2. Permítanme advertirles que no lo están. Recordarán que la alteración del barniz en el n.° 1 se manifestaba como aberturas en forma de cráter. Solo en un punto se encuentran similares estos dos fenómenos: ambos se localizan en la caja de resonancia. A mi parecer, esto constituye una prueba más del papel fundamental que desempeña la caja de resonancia del violín en la producción del sonido. Ambas alteraciones del barniz apuntan a una mayor amplitud y potencia de oscilación de la caja de resonancia como causa de su existencia.
El fenómeno de barniz n.° 2 se presenta en forma de gotas colgantes, que cuelgan de la superficie interior de la caja de resonancia como estalactitas de la cúpula de la Cueva Mammoth. Estas gotas no se encuentran en toda la caja de resonancia. Solo existen en el área central y dentro de las líneas trazadas sobre el centro de las salidas. La gradación de esta caja de resonancia...
121

¡El tablero de instrumentos! es de 2 y 8-10 mm a 2 y 7-10. Por qué estos fenómenos tv/6 varían tan diametralmente en apariencia física está más allá de mi comprensión.
Desde este comienzo rudimentario continué hasta sentir que la derrota se convertía en victoria. En la etapa experimental de este trabajo probé varias sustancias para la protección de la superficie interior. Una de estas sustancias, o combinación de dos, aunque descartada por considerarse inútil, aún posee interés porque muestra el efecto sobre el tono del violín al doblar la caja de resonancia. Esta mezcla está compuesta de aceite de linaza hervido y cola transparente. La unión de estas dos sustancias es puramente mecánica. La presencia de aceite provoca flexibilidad en la masa cuando está seca. Tal es la flexibilidad que, cuando se secan en láminas delgadas, pueden doblarse indefinidamente sobre sí mismas sin romperse. Debido a esta cualidad física, probé esta mezcla extensamente. En esta combinación, la cantidad proporcional es de 25 gotas de aceite por una onza líquida de cola, con la consistencia adecuada para ser aplicada con un pincel.
Hay dos objeciones a esta mezcla como material para la protección de la superficie interior del violín. Cualquiera de estas objeciones es suficiente para condenar su uso en el violín. Primero: La humedad del pegamento es absorbida por la madera. Segundo: Al secarse, esta mezcla se contrae con la fuerza suficiente para doblar la tapa armónica. Mediante capas repetidas, las "concavidades" de la tapa armónica pueden profundizarse hasta que el arqueamiento se eleve 1/4 de pulgada. Con la
) la caja de resonancia así doblada, el tono sufre- un
122

un cambio sorprendente. Estos cambios se manifiestan por una corta duración y por una cualidad estridente (como la voz ronca) y por un mayor brillo o distinción en el pasaje tocado rápidamente.
Hay violinistas que se complacen con la corta duración del sonido que produce la caja de resonancia curvada. Estos intérpretes se distinguen por su rapidez de ejecución. Como es bien sabido, existen violines que solo producen una mezcla confusa de sonidos en los pasajes rápidos. Comparando estos violines con aquellos que poseen un sonido estridente pero brillante, prefiero estos últimos. Sin embargo, en mi opinión, el sonido estridente y brillante no se compara en absoluto con el sonido natural y brillante. El sonido estridente es débil y, por lo tanto, artificial.
Por alguna razón misteriosa, los distintos gustos tonales del violín se aproximan más al infinito que los gustos tonales de cualquier otro instrumento musical.
Mi amigo JD O'Brien, luthier aficionado de Pittsburg, Pensilvania, tuvo durante su último viaje europeo el placer, o más bien el dolor, de ver cinco Stradivarius desgastados al cuidado de Hammeg, luthier de Berlín. Uno de estos violines fue obsequiado en una ocasión al Dr. Joachim por sus admiradores londinenses.
En los últimos años he preguntado: 4 '¿Por qué nunca leemos sobre el desgastado Amati, o el desgastado Guaneri, o el desgastado Maggini, o incluso el desgastado Da Salo?' No recibo respuesta a ello. La suposición parece ser la única fuente de respuesta. Como ya se ha dicho, esa caja de resonancia que produce el tono más agradable es la primera en sucumbir a las fuerzas desintegradoras del calor y
123

Humedad. Por supuesto, el violín que produce el tono más agradable probablemente recibirá mayor uso en un período de tiempo determinado. Pero el uso legítimo por sí solo no puede explicar la rapidez de su deterioro. Si me llamara Hammeg, antes de cumplir una hora, esas cajas de resonancia desgastadas serían sometidas a pruebas de densidad de fibra dura y tejido conectivo entre ellas.
Hoy estoy convencido de que toda caja de resonancia desgastada, si se toma a tiempo y se le da la protección adecuada de la superficie interior, podría haber conservado todo su valor tonal original por un período indefinido. Sobre este punto, he hecho una amplia demostración práctica y me he asegurado de que tal conclusión se basa en razones seguras. He dado detalles correctos y erróneos sobre la protección de la superficie interior para que usted pueda comprobarlo con mucha menos dificultad. Puede preguntarse: "¿Qué lógica hay en proteger una superficie de las tapas del violín mientras se deja otra superficie sin ninguna protección?". Recibirá varias respuestas a esta pregunta. Con toda probabilidad, recibirá respuestas positivas de personas que prácticamente no saben nada sobre este tema de la protección de la superficie interior de los violines, que dicha protección atrapa las secreciones en la madera. Ciertamente las atrapa, es decir, si hay alguna secreción en ella. Si me ha leído atentamente, habrá observado que tengo cuidado de repetir el hecho de que solo aplico protección de la superficie interior a violines usados, violines que han tenido tiempo de completar el proceso de contracción después de salir de las manos del constructor; y
124
\label{barniz_cremona}
que tengo cuidado de secar aún más la madera con calor artificial cuando sea necesario.
Si realizas el trabajo de protección de la superficie interior* siguiendo cuidadosamente las instrucciones, te llevarás una grata sorpresa. Tu viejo violín adquirirá una mayor intensidad y brillo en el sonido; y si su timbre ya era valioso antes de proteger y alisar las superficies interiores de forma perfecta y permanente, disfrutarás sabiendo que esta mejora en el tono es permanente.
La permanencia del timbre del violín es una cualidad muy deseable. Está a tu alcance si la pides.
Es posible que no compartan mis sentimientos respecto a la insensatez de ciertas afirmaciones sobre el barniz de violín Cremona. Como estudiantes de violín, este tema les ha sido mencionado muchas veces. Han leído sobre él en libros, folletos, revistas e incluso en entrevistas de periódicos. Es posible que dicha lectura no les haya causado la misma impresión que a mí. Nada en la vida es más evidente que el hecho de que las impresiones mentales pueden variar según el número de mentes. En mi juventud leí algo sobre la belleza del barniz de violín Cremona. Entonces pensé que el autor se refería realmente a las gomas que componen el barniz. En la mediana edad, comencé a leer que la superioridad del tono del violín Cremona se debe al barniz Cremona. En la edad adulta, comencé a leer que el barniz Cremona, de alguna manera, penetraba la madera; también, que el barniz de violín Cremona...
125

Desapareció misteriosamente de los canales comerciales. Curiosamente, no he leído nada sobre la desaparición de otros barnices para muebles. ¡Es sumamente extraño que el destino haya elegido el barniz para violines para su aniquilación! Más allá de toda extrañeza está el hecho de que solo un número limitado de violines de Cremona son, o fueron alguna vez, célebres por la superioridad de su sonido. Señalando directamente a los dos constructores de violines cremoneses más famosos, Stradivarius y Joseph Guarnerius, "expertos" de confianza me han dicho: "No todos los violines tenían el mismo valor tonal".
"¿Por qué no?"
¡Tenían acceso al barniz de Cremona!
Realmente~ahora podemos esperar una "entrevista" de algún "experto" londinense (más bien promotor comercial), en la que afirme que el barniz de Cremona era tan caro (¡oh!) que nadie podía permitirse usarlo excepto de forma limitada. . ,
Quizás tenga razón.
El uso de un barniz que penetra en la madera requiere sin duda una gran cantidad de este.
Los italianos son reconocidos, y lo han sido durante mucho tiempo, por su superior sentido del color. Sin embargo, incluso en el trabajo de color en violines, no todos los violines eran igualmente bellos. No me considero un artista del color, pero, por mi experiencia en el trabajo de color sobre madera, estoy seguro de que nadie puede hacer lo que se llama un violín hermoso sin tener una madera hermosa como requisito previo. Incluso así, es fácil fracasar. Seis veces he tenido que quitar el trabajo de color de un solo violín antes de encontrar el adecuado.
126

combinación para que coincida con esa muestra de madera en particular; y puede transcurrir un largo intervalo antes de que esa combinación sea igualmente efectiva en otra muestra. La gente exclama: "¡Qué barniz tan bonito!"
La belleza no reside solo en las encías. Está en la combinación de colores.
Ni un ápice restaría mérito alguno a los "viejos maestros". Tampoco puedo atribuirles un mérito inmerecido. Por esas declaraciones vacías, cuyo único propósito es la "promoción comercial" de los violines de Cremona, solo siento desprecio.
Sin duda, mi desprecio por esas "entrevistas" londinenses y por esos autores londinenses cuyo animosidad es claramente el "comercio" es comparable al de ellos por los estadounidenses. La opinión que los londinenses tienen de los estadounidenses queda claramente reflejada en las siguientes palabras atribuidas a Kipling: «Estaba a bordo de un velero que viajaba de Calcuta a Londres. Una mañana, frente a la costa occidental de África, la guardia informó haber sido pasada al amanecer por la serpiente marina más grande jamás registrada. La guardia declaró que su atención fue atraída por primera vez al monstruo por un silbido parecido al vapor cuando emergió bajo la popa. Su lengua bífida medía ocho pies de largo. Su cabeza medía diez pies de largo desde la punta de la mandíbula hasta el único ojo llameante. Su cabeza se elevaba quince pies sobre el agua. Su cuerpo era más largo que el barco. Mientras escuchaba esta historia de marineros, un periodista neoyorquino, ocupado con lápiz y cuaderno, le comentó a Kipling:
127

que al llegar a Londres informaría de este extraordinario incidente a la prensa. —No lo hagas —dijo Kipling. —¿Por qué no? —replicó el neoyorquino. —Amigo mío —respondió Kipling—, déjame contarte algo que parece que desconoces. Inglaterra ya tenía canas antes de que nacieras. Guarda esta historia para cuando llegues a casa.
Debemos admitir la veracidad de la observación de Kipling. Inglaterra ostenta el monopolio de las canas. No será culpa suya si alguna vez aparece una sola cana en otro país. Hace más de cien años, recibió una insinuación de los estadounidenses de que deseábamos vivir nuestras vidas con naturalidad, independientemente de si esa vida sería lo suficientemente larga como para llegar a la etapa de las canas. A juzgar por las señales del presente, hay esperanza para nosotros. Desde Portland hasta Seattle, desde Nueva Orleans hasta Duluth, hay quienes ya no se creen la historia del experto londinense sobre la "impregnación de la madera por el barniz". Hoy, el experto londinense, "entrevistándose" a sí mismo, se lamenta (!) de los precios prohibitivos de los pocos violines antiguos impregnados por el barniz Cremona perdido. A juzgar por el pasado, ¿qué daño hay en profetizar el futuro? Convicto, predigo con seguridad que, en el presente siglo, los europeos buscarán el mejor violín por toda América.
¿Por qué no?
Los violines de Cremona, con barniz aplicado solo en una superficie, pronto estarán fuera del mercado. Es seguro que algunos violines americanos, con barniz aplicado en ambas superficies, estarán entonces en el mercado. ¿Por qué no imaginar de-
128

¿Los descendientes del actual experto londinense se entrevistan entonces a sí mismos sobre los precios prohibitivos de los violines estadounidenses de doble perforación? Hay un estadounidense de pelo canoso que disfrutaría enormemente de esa imagen.
